%%%%%%%%%%%%%%%%%%%%%%%%%%%%%%%%%%%%%%%%%%%%%%%%%%%%%%%%%%%%
\chapter{Introduction}
\label{chap:introduction}

%%%%%%%%%%%%%%%%%%%%%%%%%%%%%%%%%%%%%%%%%%%%%%%%%%%%%%%%%%%%
\section{Purpose}
\label{sec:srs-purpose}

The purpose of this Software Requirements Specification is to define
the functional, non-functional, interface, database, security, and
operational requirements for the EcoTrace Intelligent Electronic Waste
Recovery and Circular Economy Management System.

This document provides a common reference for developers, testers,
supervisors, system administrators, project stakeholders, and future
maintainers.

%%%%%%%%%%%%%%%%%%%%%%%%%%%%%%%%%%%%%%%%%%%%%%%%%%%%%%%%%%%%
\section{Scope}
\label{sec:srs-scope}

EcoTrace is an intelligent cross-platform system designed to support
electronic waste registration, collection, transportation, repair,
refurbishment, recycling, resale, reporting, and circular economy
operations.

The system supports household users, businesses, collectors, drivers,
technicians, recyclers, administrators, and external organizations.

%%%%%%%%%%%%%%%%%%%%%%%%%%%%%%%%%%%%%%%%%%%%%%%%%%%%%%%%%%%%
\section{Intended Audience}
\label{sec:intended-audience}

This document is intended for:

\begin{itemize}

\item software developers;

\item software testers;

\item system analysts;

\item project supervisors;

\item database administrators;

\item environmental agencies;

\item system administrators;

\item project stakeholders; and

\item future system maintainers.

\end{itemize}

%%%%%%%%%%%%%%%%%%%%%%%%%%%%%%%%%%%%%%%%%%%%%%%%%%%%%%%%%%%%
\section{Definitions}
\label{sec:definitions}

\begin{table}[H]
\centering
\caption{Definitions}

\begin{tabularx}{\textwidth}{
    |p{4cm}|
    X|
}

\hline
\textbf{Term} &
\textbf{Definition} \\

\hline
Electronic Waste &
Discarded electrical or electronic equipment and components. \\

\hline
Circular Economy &
An economic approach that promotes reuse, repair, refurbishment,
recycling, and responsible resource recovery. \\

\hline
Reverse Logistics &
The movement of used products from consumers back to collection,
repair, refurbishment, or recycling facilities. \\

\hline
Household User &
A residential user participating in EcoTrace collection and recycling
activities. \\

\hline
Collector &
An authorized field worker responsible for collecting electronic
waste. \\

\hline

\end{tabularx}
\end{table}

%%%%%%%%%%%%%%%%%%%%%%%%%%%%%%%%%%%%%%%%%%%%%%%%%%%%%%%%%%%%
\section{Acronyms and Abbreviations}
\label{sec:acronyms}

\begin{table}[H]
\centering
\caption{Acronyms and Abbreviations}

\begin{tabularx}{\textwidth}{
    |p{3cm}|
    X|
}

\hline
\textbf{Acronym} &
\textbf{Meaning} \\

\hline
AI & Artificial Intelligence \\

\hline
API & Application Programming Interface \\

\hline
CRUD & Create, Read, Update, and Delete \\

\hline
FCM & Firebase Cloud Messaging \\

\hline
GIS & Geographic Information System \\

\hline
GPS & Global Positioning System \\

\hline
IoT & Internet of Things \\

\hline
RBAC & Role-Based Access Control \\

\hline
SRS & Software Requirements Specification \\

\hline
UML & Unified Modeling Language \\

\hline
UI & User Interface \\

\hline
UX & User Experience \\

\hline

\end{tabularx}
\end{table}

%%%%%%%%%%%%%%%%%%%%%%%%%%%%%%%%%%%%%%%%%%%%%%%%%%%%%%%%%%%%
\section{References}
\label{sec:srs-references}

This SRS is informed by the following standards and technical
documentation:

\begin{itemize}

\item IEEE 29148: Systems and Software Engineering Requirements
Engineering;

\item ISO/IEC 25010 Software Product Quality Model;

\item Flutter technical documentation;

\item Firebase technical documentation;

\item PyTorch technical documentation;

\item Google Maps Platform documentation; and

\item Njala University project-writing guidelines.

\end{itemize}

%%%%%%%%%%%%%%%%%%%%%%%%%%%%%%%%%%%%%%%%%%%%%%%%%%%%%%%%%%%%
\section{Document Overview}
\label{sec:document-overview}

This SRS is organized into nine chapters.

Chapter 1 introduces the document, scope, definitions, references, and
intended audience.

Chapter 2 presents the overall product description.

Chapter 3 defines the external interface requirements.

Chapter 4 specifies the system's functional requirements.

Chapter 5 presents the non-functional requirements.

Chapter 6 defines the database requirements.

Chapter 7 specifies the business rules.

Chapter 8 presents detailed use case specifications.

Chapter 9 contains supporting appendices and traceability information.

%%%%%%%%%%%%%%%%%%%%%%%%%%%%%%%%%%%%%%%%%%%%%%%%%%%%%%%%%%%%